\IfFileExists{stacks-project.cls}{%
\documentclass{stacks-project}
}{%
\documentclass{amsart}
}

% For dealing with references we use the comment environment
\usepackage{verbatim}
\newenvironment{reference}{\comment}{\endcomment}
%\newenvironment{reference}{}{}
\newenvironment{slogan}{\comment}{\endcomment}
\newenvironment{history}{\comment}{\endcomment}

% For commutative diagrams we use Xy-pic
\usepackage[all]{xy}

% We use 2cell for 2-commutative diagrams.
\xyoption{2cell}
\UseAllTwocells

% We use multicol for the list of chapters between chapters
\usepackage{multicol}

% This is generally recommended for better output
\usepackage{lmodern}
\usepackage[T1]{fontenc}

% For cross-file-references
\usepackage{xr-hyper}

% Package for hypertext links:
\usepackage{hyperref}

% For any local file, say "hello.tex" you want to link to please
% use \externaldocument[hello-]{hello}
\externaldocument[introduction-]{introduction}
\externaldocument[conventions-]{conventions}
\externaldocument[sets-]{sets}
\externaldocument[categories-]{categories}
\externaldocument[topology-]{topology}
\externaldocument[sheaves-]{sheaves}
\externaldocument[sites-]{sites}
\externaldocument[stacks-]{stacks}
\externaldocument[fields-]{fields}
\externaldocument[algebra-]{algebra}
\externaldocument[brauer-]{brauer}
\externaldocument[homology-]{homology}
\externaldocument[derived-]{derived}
\externaldocument[simplicial-]{simplicial}
\externaldocument[more-algebra-]{more-algebra}
\externaldocument[smoothing-]{smoothing}
\externaldocument[modules-]{modules}
\externaldocument[sites-modules-]{sites-modules}
\externaldocument[injectives-]{injectives}
\externaldocument[cohomology-]{cohomology}
\externaldocument[sites-cohomology-]{sites-cohomology}
\externaldocument[dga-]{dga}
\externaldocument[dpa-]{dpa}
\externaldocument[sdga-]{sdga}
\externaldocument[hypercovering-]{hypercovering}
\externaldocument[schemes-]{schemes}
\externaldocument[constructions-]{constructions}
\externaldocument[properties-]{properties}
\externaldocument[morphisms-]{morphisms}
\externaldocument[coherent-]{coherent}
\externaldocument[divisors-]{divisors}
\externaldocument[limits-]{limits}
\externaldocument[varieties-]{varieties}
\externaldocument[topologies-]{topologies}
\externaldocument[descent-]{descent}
\externaldocument[perfect-]{perfect}
\externaldocument[more-morphisms-]{more-morphisms}
\externaldocument[flat-]{flat}
\externaldocument[groupoids-]{groupoids}
\externaldocument[more-groupoids-]{more-groupoids}
\externaldocument[etale-]{etale}
\externaldocument[chow-]{chow}
\externaldocument[intersection-]{intersection}
\externaldocument[pic-]{pic}
\externaldocument[weil-]{weil}
\externaldocument[adequate-]{adequate}
\externaldocument[dualizing-]{dualizing}
\externaldocument[duality-]{duality}
\externaldocument[discriminant-]{discriminant}
\externaldocument[derham-]{derham}
\externaldocument[local-cohomology-]{local-cohomology}
\externaldocument[algebraization-]{algebraization}
\externaldocument[curves-]{curves}
\externaldocument[resolve-]{resolve}
\externaldocument[models-]{models}
\externaldocument[functors-]{functors}
\externaldocument[equiv-]{equiv}
\externaldocument[pione-]{pione}
\externaldocument[etale-cohomology-]{etale-cohomology}
\externaldocument[proetale-]{proetale}
\externaldocument[relative-cycles-]{relative-cycles}
\externaldocument[more-etale-]{more-etale}
\externaldocument[trace-]{trace}
\externaldocument[crystalline-]{crystalline}
\externaldocument[spaces-]{spaces}
\externaldocument[spaces-properties-]{spaces-properties}
\externaldocument[spaces-morphisms-]{spaces-morphisms}
\externaldocument[decent-spaces-]{decent-spaces}
\externaldocument[spaces-cohomology-]{spaces-cohomology}
\externaldocument[spaces-limits-]{spaces-limits}
\externaldocument[spaces-divisors-]{spaces-divisors}
\externaldocument[spaces-over-fields-]{spaces-over-fields}
\externaldocument[spaces-topologies-]{spaces-topologies}
\externaldocument[spaces-descent-]{spaces-descent}
\externaldocument[spaces-perfect-]{spaces-perfect}
\externaldocument[spaces-more-morphisms-]{spaces-more-morphisms}
\externaldocument[spaces-flat-]{spaces-flat}
\externaldocument[spaces-groupoids-]{spaces-groupoids}
\externaldocument[spaces-more-groupoids-]{spaces-more-groupoids}
\externaldocument[bootstrap-]{bootstrap}
\externaldocument[spaces-pushouts-]{spaces-pushouts}
\externaldocument[spaces-chow-]{spaces-chow}
\externaldocument[groupoids-quotients-]{groupoids-quotients}
\externaldocument[spaces-more-cohomology-]{spaces-more-cohomology}
\externaldocument[spaces-simplicial-]{spaces-simplicial}
\externaldocument[spaces-duality-]{spaces-duality}
\externaldocument[formal-spaces-]{formal-spaces}
\externaldocument[restricted-]{restricted}
\externaldocument[spaces-resolve-]{spaces-resolve}
\externaldocument[formal-defos-]{formal-defos}
\externaldocument[defos-]{defos}
\externaldocument[cotangent-]{cotangent}
\externaldocument[examples-defos-]{examples-defos}
\externaldocument[algebraic-]{algebraic}
\externaldocument[examples-stacks-]{examples-stacks}
\externaldocument[stacks-sheaves-]{stacks-sheaves}
\externaldocument[criteria-]{criteria}
\externaldocument[artin-]{artin}
\externaldocument[quot-]{quot}
\externaldocument[stacks-properties-]{stacks-properties}
\externaldocument[stacks-morphisms-]{stacks-morphisms}
\externaldocument[stacks-limits-]{stacks-limits}
\externaldocument[stacks-cohomology-]{stacks-cohomology}
\externaldocument[stacks-perfect-]{stacks-perfect}
\externaldocument[stacks-introduction-]{stacks-introduction}
\externaldocument[stacks-more-morphisms-]{stacks-more-morphisms}
\externaldocument[stacks-geometry-]{stacks-geometry}
\externaldocument[moduli-]{moduli}
\externaldocument[moduli-curves-]{moduli-curves}
\externaldocument[examples-]{examples}
\externaldocument[exercises-]{exercises}
\externaldocument[guide-]{guide}
\externaldocument[desirables-]{desirables}
\externaldocument[coding-]{coding}
\externaldocument[obsolete-]{obsolete}
\externaldocument[fdl-]{fdl}
\externaldocument[index-]{index}

% Theorem environments.
%
\theoremstyle{plain}
\newtheorem{theorem}[subsection]{Theorem}
\newtheorem{proposition}[subsection]{Proposition}
\newtheorem{lemma}[subsection]{Lemma}

\theoremstyle{definition}
\newtheorem{definition}[subsection]{Definition}
\newtheorem{example}[subsection]{Example}
\newtheorem{exercise}[subsection]{Exercise}
\newtheorem{situation}[subsection]{Situation}

\theoremstyle{remark}
\newtheorem{remark}[subsection]{Remark}
\newtheorem{remarks}[subsection]{Remarks}

\numberwithin{equation}{subsection}

% Macros
%
\def\lim{\mathop{\mathrm{lim}}\nolimits}
\def\colim{\mathop{\mathrm{colim}}\nolimits}
\def\Spec{\mathop{\mathrm{Spec}}}
\def\Hom{\mathop{\mathrm{Hom}}\nolimits}
\def\Ext{\mathop{\mathrm{Ext}}\nolimits}
\def\SheafHom{\mathop{\mathcal{H}\!\mathit{om}}\nolimits}
\def\SheafExt{\mathop{\mathcal{E}\!\mathit{xt}}\nolimits}
\def\Sch{\mathit{Sch}}
\def\Mor{\mathop{\mathrm{Mor}}\nolimits}
\def\Ob{\mathop{\mathrm{Ob}}\nolimits}
\def\Sh{\mathop{\mathit{Sh}}\nolimits}
\def\NL{\mathop{N\!L}\nolimits}
\def\CH{\mathop{\mathrm{CH}}\nolimits}
\def\proetale{{pro\text{-}\acute{e}tale}}
\def\etale{{\acute{e}tale}}
\def\QCoh{\mathit{QCoh}}
\def\Ker{\mathop{\mathrm{Ker}}}
\def\Im{\mathop{\mathrm{Im}}}
\def\Coker{\mathop{\mathrm{Coker}}}
\def\Coim{\mathop{\mathrm{Coim}}}

% Boxtimes
%
\DeclareMathSymbol{\boxtimes}{\mathbin}{AMSa}{"02}

%
% Macros for moduli stacks/spaces
%
\def\QCohstack{\mathcal{QC}\!\mathit{oh}}
\def\Cohstack{\mathcal{C}\!\mathit{oh}}
\def\Spacesstack{\mathcal{S}\!\mathit{paces}}
\def\Quotfunctor{\mathrm{Quot}}
\def\Hilbfunctor{\mathrm{Hilb}}
\def\Curvesstack{\mathcal{C}\!\mathit{urves}}
\def\Polarizedstack{\mathcal{P}\!\mathit{olarized}}
\def\Complexesstack{\mathcal{C}\!\mathit{omplexes}}
% \Pic is the operator that assigns to X its picard group, usage \Pic(X)
% \Picardstack_{X/B} denotes the Picard stack of X over B
% \Picardfunctor_{X/B} denotes the Picard functor of X over B
\def\Pic{\mathop{\mathrm{Pic}}\nolimits}
\def\Picardstack{\mathcal{P}\!\mathit{ic}}
\def\Picardfunctor{\mathrm{Pic}}
\def\Deformationcategory{\mathcal{D}\!\mathit{ef}}
\usepackage{cleveref}
\crefname{originallemma}{original lemma}{original lemmas}
\Crefname{originallemma}{Original lemma}{Original lemmas}
\makeatletter
\newlabel{lemma-separate}{{0CM8}{source}{}{sourcepointer-1}{}}
\newlabel{lemma-separate@cref}{{[originallemma][0][]0CM8}{[1][0][]source}{}{}{}}
\makeatother
\crefname{originalreference}{original reference}{original references}
\Crefname{originalreference}{Original reference}{Original references}
\makeatletter
\newlabel{morphisms-lemma-quasi-affine-finite-type-over-S}{{04II}{source}{}{sourcepointer-2}{}}
\newlabel{morphisms-lemma-quasi-affine-finite-type-over-S@cref}{{[originalreference][0][]04II}{[1][0][]source}{}{}{}}
\makeatother
\crefname{originalreference}{original reference}{original references}
\Crefname{originalreference}{Original reference}{Original references}
\makeatletter
\newlabel{spaces-morphisms-lemma-universally-closed-local}{{03IT}{source}{}{sourcepointer-3}{}}
\newlabel{spaces-morphisms-lemma-universally-closed-local@cref}{{[originalreference][0][]03IT}{[1][0][]source}{}{}{}}
\makeatother
\crefname{originalreference}{original reference}{original references}
\Crefname{originalreference}{Original reference}{Original references}
\makeatletter
\newlabel{spaces-properties-lemma-points-cartesian}{{03H4}{source}{}{sourcepointer-4}{}}
\newlabel{spaces-properties-lemma-points-cartesian@cref}{{[originalreference][0][]03H4}{[1][0][]source}{}{}{}}
\makeatother
\crefname{originalreference}{original reference}{original references}
\Crefname{originalreference}{Original reference}{Original references}
\makeatletter
\newlabel{spaces-properties-lemma-subspace-induced-topology}{{0ABJ}{source}{}{sourcepointer-5}{}}
\newlabel{spaces-properties-lemma-subspace-induced-topology@cref}{{[originalreference][0][]0ABJ}{[1][0][]source}{}{}{}}
\makeatother
\title{Testing universal closedness on affine space fibres}
\author{The Stacks Project Authors}
\date{stacks-project:spaces-limits.tex:local-snapshot-20261012; GNU FDL 1.2 or later}
\begin{document}
\maketitle
\noindent\textit{Editorial note.} Counted span is the author's own proof environment, source lines 4643-4732, following the statement at 4628-4641.\par\medskip
\noindent\textit{Editorial note.} The '(some details omitted)' covers the etale-local-on-Z reduction of closedness of the topological map; flagged in the notes.\par\medskip
\noindent\textit{Editorial note.} Steps the author leaves to the reader are flagged here for the reviewer: source line(s) 4699 carry an "omitted"/"follows easily"/"immediate from the definitions" phrase. Those sentences are retained verbatim; the displayed argument around them stays inside the counted span.\par\medskip
\noindent\textit{Allowed standard prerequisites (editorial).}\begin{itemize}
\item Separation lemma for non-closed images. lemma-separate (same chapter, cited at \textasciitilde{}4650-4651). Produces the etale neighbourhood (V,v), the diagram through Z', and T'.
\item Points of cartesian squares. Properties of Spaces, Lemma spaces-properties-lemma-points-cartesian (cross-chapter, cited at \textasciitilde{}4679-4680). Identifies the image of T\_V in |V| with the pullback of the image of T.
\item Universal closedness is etale local on the target. Morphisms of Spaces, Lemma spaces-morphisms-lemma-universally-closed-local (cross-chapter, cited at \textasciitilde{}4690-4692). Lifts universal closedness of X x\_Y V -> V back to f.
\item Quasi-affine immersion into affine space over the base. Morphisms, Lemma morphisms-lemma-quasi-affine-finite-type-over-S (cross-chapter, cited at \textasciitilde{}4701-4702). Immerses Z into A\textasciicircum{}n x Y.
\item Subspace topology on topological images. Properties of Spaces, Lemma spaces-properties-lemma-subspace-induced-topology (cross-chapter, cited at \textasciitilde{}4725-4726). Horizontal arrows are homeomorphisms onto locally closed subsets.
\end{itemize}\medskip
\ifdefined\hypertarget\hypertarget{sourcepointer-1}{}\fi
\noindent\textit{Source reference.} Original lemma 0CM8: sources/stacks-spaces-limits/pinned-original-source.tex, source lines 4529--4529. Navigation pointer only: the counted proof cites this statement by \textbackslash{}ref as a standing input; the deduction the counted proof has to make is carried out inside the counted span, and no part of that deduction is delegated to this pointer.\par\medskip
\ifdefined\hypertarget\hypertarget{sourcepointer-2}{}\fi
\noindent\textit{Source reference.} Original reference 04II: focused-seeds/geometry/S9-spaces-limits-test-universally-closed/ref-tags.txt, source lines 5--5. Navigation pointer only: the counted proof cites this statement by \textbackslash{}ref as a standing input; the deduction the counted proof has to make is carried out inside the counted span, and no part of that deduction is delegated to this pointer.\par\medskip
\ifdefined\hypertarget\hypertarget{sourcepointer-3}{}\fi
\noindent\textit{Source reference.} Original reference 03IT: focused-seeds/geometry/S9-spaces-limits-test-universally-closed/ref-tags.txt, source lines 6--6. Navigation pointer only: the counted proof cites this statement by \textbackslash{}ref as a standing input; the deduction the counted proof has to make is carried out inside the counted span, and no part of that deduction is delegated to this pointer.\par\medskip
\ifdefined\hypertarget\hypertarget{sourcepointer-4}{}\fi
\noindent\textit{Source reference.} Original reference 03H4: focused-seeds/geometry/S9-spaces-limits-test-universally-closed/ref-tags.txt, source lines 7--7. Navigation pointer only: the counted proof cites this statement by \textbackslash{}ref as a standing input; the deduction the counted proof has to make is carried out inside the counted span, and no part of that deduction is delegated to this pointer.\par\medskip
\ifdefined\hypertarget\hypertarget{sourcepointer-5}{}\fi
\noindent\textit{Source reference.} Original reference 0ABJ: focused-seeds/geometry/S9-spaces-limits-test-universally-closed/ref-tags.txt, source lines 8--8. Navigation pointer only: the counted proof cites this statement by \textbackslash{}ref as a standing input; the deduction the counted proof has to make is carried out inside the counted span, and no part of that deduction is delegated to this pointer.\par\medskip
\begin{lemma}
\label{lemma-test-universally-closed}
Let $S$ be a scheme.
Let $f : X \to Y$ be a quasi-compact morphism of algebraic spaces over $S$.
The following are equivalent
\begin{enumerate}
\item $f$ is universally closed,
\item for every morphism $Z \to Y$ which is locally of finite presentation
the map $|X \times_Y Z| \to |Z|$ is closed, and
\item there exists a scheme $V$ and a surjective \'etale morphism $V \to Y$
such that $|\mathbf{A}^n \times (X \times_Y V)| \to |\mathbf{A}^n \times V|$
is closed for all $n \geq 0$.
\end{enumerate}
\end{lemma}

\begin{proof}
It is clear that (1) implies (2).
Suppose that $|X \times_Y Z| \to |Z|$ is not closed for some
morphism of algebraic spaces $Z \to Y$ over $S$. This means that there
exists some closed subset $T \subset |X \times_Y Z|$
such that $\Im(T \to |Z|)$ is not closed. Pick $z \in |Z|$
in the closure of the image of $T$ but not in the image.
Apply Lemma \ref{lemma-separate}.
We find an \'etale neighbourhood $(V, v) \to (Z, z)$, a commutative diagram
$$
\xymatrix{
V \ar[d] \ar[r]_a & Z' \ar[d]^b \\
Z \ar[r]^g & Y,
}
$$
and a closed subset $T' \subset |X \times_Y Z'|$ such that
\begin{enumerate}
\item the morphism $b : Z' \to Y$ is locally of finite presentation,
\item with $z' = a(v)$ we have $z' \not \in \Im(T' \to |Z'|)$, and
\item the inverse image of $T$ in $|X \times_Y V|$ maps into $T'$ via
$|X \times_Y V| \to |X \times_Y Z'|$.
\end{enumerate}
We claim that $z'$ is in the closure of $\Im(T' \to |Z'|)$
which implies that $|X \times_Y Z'| \to |Z'|$ is not closed.
The claim shows that (2) implies (1).
To see the claim is true we contemplate
following commutative diagram
$$
\xymatrix{
X \times_Y Z \ar[d] &
X \times_Y V \ar[l] \ar[d] \ar[r] &
X \times_Y Z' \ar[d] \\
Z & V \ar[l] \ar[r]^a & Z'
}
$$
Let $T_V \subset |X \times_Y V|$ be the inverse image of $T$.
By Properties of Spaces, Lemma \ref{spaces-properties-lemma-points-cartesian}
the image of $T_V$ in $|V|$ is the inverse image of the image
of $T$ in $|Z|$. Then since $z$ is in the closure of the image of
$T \to |Z|$ and since $|V| \to |Z|$ is open, we see that $v$ is in
the closure of the image of $T_V \to |V|$. Since the image of
$T_V$ in $|X \times_Y Z'|$ is contained in $|T'|$ it follows
immediately that $z' = a(v)$ is in the closure of the image of $T'$.

\medskip\noindent
It is clear that (1) implies (3). Let $V \to Y$ be as in (3).
If we can show that $X \times_Y V \to V$ is universally closed,
then $f$ is universally closed by
Morphisms of Spaces, Lemma
\ref{spaces-morphisms-lemma-universally-closed-local}.
Thus it suffices to show that $f : X \to Y$ satisfies (2)
if $f$ is a quasi-compact morphism of algebraic spaces,
$Y$ is a scheme, and $|\mathbf{A}^n \times X| \to |\mathbf{A}^n \times Y|$
is closed for all $n$. Let $Z \to Y$ be locally of finite presentation.
We have to show the map $|X \times_Y Z| \to |Z|$ is closed.
This question is \'etale local on $Z$ hence we may assume $Z$
is affine (some details omitted). Since $Y$ is a scheme, $Z$ is affine,
and $Z \to Y$ is locally of finite presentation we can find
an immersion $Z \to \mathbf{A}^n \times Y$, see
Morphisms, Lemma \ref{morphisms-lemma-quasi-affine-finite-type-over-S}.
Consider the cartesian diagram
$$
\vcenter{
\xymatrix{
X \times_Y Z \ar[d] \ar[r] & \mathbf{A}^n \times X \ar[d] \\
Z \ar[r] & \mathbf{A}^n \times Y
}
}
\quad
\begin{matrix}
\text{inducing the} \\
\text{cartesian square}
\end{matrix}
\quad
\vcenter{
\xymatrix{
|X \times_Y Z| \ar[d] \ar[r] & |\mathbf{A}^n \times X| \ar[d] \\
|Z| \ar[r] & |\mathbf{A}^n \times Y|
}
}
$$
of topological spaces whose horizontal arrows are homeomorphisms
onto locally closed subsets (Properties of Spaces, Lemma
\ref{spaces-properties-lemma-subspace-induced-topology}).
Thus every closed subset $T$
of $|X \times_Y Z|$ is the pullback of a closed subset $T'$ of
$|\mathbf{A}^n \times Y|$. Since the assumption is that the image
of $T'$ in $|\mathbf{A}^n \times X|$ is closed we conclude that
the image of $T$ in $|Z|$ is closed as desired.
\end{proof}

\end{document}
